% titel: Blatt 2, A1: Summenformeln aus Ionen
% Teil II: der durchgerechnete Weg.

\subsection*{Gegeben und gesucht}
Kationen und Anionen; gesucht jeweils die ladungsneutrale Summenformel.

\subsection*{Rechnung}
Ladungsneutralitaet uber einfachen Ladungsausgleich:
\begin{itemize}
  \item[a)] $\mathrm{Fe^{3+}}$ + $\mathrm{Cl^-}$: ein Eisen (3+) braucht drei
        Chloride (je 1-). $\Rightarrow \mathrm{FeCl_3}$.
  \item[b)] $\mathrm{Al^{3+}}$ + $\mathrm{O^{2-}}$: kleinster Ausgleich
        $2\cdot(3+) + 3\cdot(2-) = 0 \Rightarrow \mathrm{Al_2O_3}$.
  \item[c)] $\mathrm{Al^{3+}}$ + $\mathrm{PO_4^{3-}}$: gleiche Ladung
        $1:1 \Rightarrow \mathrm{AlPO_4}$.
  \item[d)] $\mathrm{Mg^{2+}}$ + $\mathrm{NO_3^-}$: zwei Nitrate pro Magnesium
        $\Rightarrow \mathrm{Mg(NO_3)_2}$.
  \item[e)] $\mathrm{Fe^{3+}}$ + $\mathrm{SO_4^{2-}}$: kleinster Ausgleich
        $2\cdot\mathrm{Fe^{3+}}$ mit $3\cdot\mathrm{SO_4^{2-}}$
        $\Rightarrow \mathrm{Fe_2(SO_4)_3}$.
  \item[f)] $\mathrm{NH_4^+}$ + $\mathrm{SO_4^{2-}}$: zwei Ammonium pro Sulfat
        $\Rightarrow \mathrm{(NH_4)_2SO_4}$.
\end{itemize}

\begin{ergebnis}
\begin{align*}
  \mathrm{a)\ FeCl_3,} &\qquad \mathrm{b)\ Al_2O_3,}\\
  \mathrm{c)\ AlPO_4,} &\qquad \mathrm{d)\ Mg(NO_3)_2,}\\
  \mathrm{e)\ Fe_2(SO_4)_3,} &\qquad \mathrm{f)\ (NH_4)_2SO_4.}
\end{align*}
\end{ergebnis}

\begin{kontrolle}
Jede Formel von rechts prüfen: Summe aller positiven und negativen Ladungen
ergibt null. Beispiel $\mathrm{Fe_2(SO_4)_3}$:
$2\cdot(+3) = 6+$, $3\cdot(-2) = 6-$, Summe $0$. Alle Formeln sind in der
kleinsten ganzen Stöchiometrie geschrieben.
\end{kontrolle}

\begin{merksatz}
Summenformel = kleinste ganze Zahlen, die die Ionenladungen exakt
(Ladungsausgleich) zueinander ausbalancieren. Mehratomige Ionen bei Vielfachheit
in Klammern schreiben, danach die Mengenzahl.
\end{merksatz}
